% -----------------------------------------------------------
\section{The Initiatory}
% -----------------------------------------------------------
	\label{TEMPLE-ENDOWMENT-INIT}

	% ----------------------
	\subsection{Initiatory-like ordinances in the Bible}
	% ----------------------
	
		\hyperref[EXODUS-29-4]{Exodus 29:4--9}:
		
			\begin{quote}
				And Aaron and his sons thou shalt bring unto the door of the tabernacle of the congregation, and shalt wash them with water. 5 And thou shalt take the garments, and put upon Aaron the coat, and the robe of the ephod, and the ephod, and the breastplate, and gird him with the curious girdle of the ephod: 6 And thou shalt put the mitre upon his head, and put the holy crown upon the mitre. 7 Then shalt thou take the anointing oil, and pour it upon his head, and anoint him. 8 And thou shalt bring his sons, and put coats upon them. 9 And thou shalt gird them with girdles, Aaron and his sons, and put the bonnets on them: and the priest's office shall be theirs for a perpetual statute: and thou shalt consecrate Aaron and his sons.
			\end{quote}
			
		\hyperref[MATTHEW-6-16]{Matthew 6:16--18}:
		
			\begin{quote}
				Moreover when ye fast, be not, as the hypocrites, of a sad countenance: for they disfigure their faces, that they may appear unto men to fast. Verily I say unto you, They have their reward 17 But thou, when thou fastest, anoint thine head, and wash thy face; 18 That thou appear not unto men to fast, but unto thy Father which is in secret: and thy Father, which seeth in secret, shall reward thee openly.
			\end{quote}
			
		\hyperref[MATTHEW-26-6]{Matthew 26:6--13}:
		
			\begin{quote}
				Now when Jesus was in Bethany, in the house of Simon the leper, 7 There came unto him a woman having an alabaster box of very precious aointment, and poured it on his head, as he sat at meat. 8 But when his disciples saw it, they had indignation, saying, To what purpose is this awaste? 9 For this ointment might have been sold for much, and given to the poor. 10 When Jesus understood it, he said unto them, Why trouble ye the woman? for she hath awrought a good work upon me. 11 For ye have the poor always with you; but me ye have not always. 12 For in that she hath poured this ointment on my body, she did it afor my burial. 13 Verily I say unto you, Wheresoever this gospel shall be preached in the whole world, there shall also this, that this woman hath done, be told for a memorial of her. 
			\end{quote}
			
		\hyperref[JOHN-13-4]{John 13:4--17}:
		
			\begin{quote}
				He riseth from supper, and laid aside his garments; and took a towel, and girded himself.
				5  After that he poureth water into a bason, and began to wash the disciples' feet, and to wipe them with the towel wherewith he was girded.
				6  Then cometh he to Simon Peter: and Peter saith unto him, Lord, dost thou wash my feet?
				7  Jesus answered and said unto him, What I do thou knowest not now; but thou shalt know hereafter.
				8  Peter saith unto him, Thou shalt never wash my feet.  Jesus answered him, If I wash thee not, thou hast no part with me.
				9  Simon Peter saith unto him, Lord, not my feet only, but also my hands and my head.
				10  Jesus saith to him, He that is washed needeth not save to wash his feet, but is clean every whit: and ye are clean, but not all.
				11  For he knew who should betray him; therefore said he, Ye are not all clean.
				12  So after he had washed their feet, and had taken his garments, and was set down again, he said unto them, Know ye what I have done to you?
				13  Ye call me Master and Lord: and ye say well; for so I am.
				14  If I then, your Lord and Master, have washed your feet; ye also ought to wash one another's feet.
				15  For I have given you an example, that ye should do as I have done to you.
				16  Verily, verily, I say unto you, The servant is not greater than his lord; neither he that is sent greater than he that sent him.
				17  If ye know these things, happy are ye if ye do them.
			\end{quote}
			
		\hyperref[REVELATION-1-5]{Revelation 1:5--6}:
		
			\begin{quote}
				And from Jesus Christ, who is the faithful witness, and the first begotten of the dead, and the prince of the kings of the earth. Unto him that loved us, and washed us from our sins in his own blood, 6 And hath made us kings and priests unto God and his Father; to him be glory and dominion for ever and ever. Amen.
			\end{quote}

	% ----------------------
	\subsection{Initiatory ceremony prior to 2005}
	% ----------------------
		\label{TEMPLE-ENDOWMENT-INIT-PRE-2005}
	
		(I am taking the text below from \href{http://www.ldsendowment.org}{ldsendowment.org}.)
	
		% -----
		\subsubsection{Washing}
		% -----
			\label{TEMPLE-ENDOWMENT-INIT-PRE-2005-WASH}
		
			\footnote{%
				\label{TEMPLE-ENDOWMENT-INIT-PRE-2005-WASH-nude}%
				A note from the website where I got the ordinance: ``Originally, initiates were presented to the washing room nude. Sometime in the early twentieth century, it became customary for initiates to wear a `shield,' a sleeveless robe open at the sides. The shield preserved modesty while allowing officiators to touch the various parts of the body named in the rite. In a slightly complicated procedure, the temple garment was placed on the initiate under the shield following the washing and anointing. Beginning in 2005, initates were instructed to clothe themselves in the garment in the privacy of their locker, \textit{before} being presented at the washing room. This means that they are clothed throughout the initiatory. As before 2005, initiates wear a shield during the initiatory, but this is now worn over the garment.''
				}%
			\textit{[An officiator places water on the initiate's head while pronouncing the following words.]]}
			
			Brother \rule{1cm}{0.15mm}, having authority,%
				\footnote{%
					\label{TEMPLE-ENDOWMENT-INIT-PRE-2005-WASH-authority}%
					Another note from the website where I got the ordinance: ``The initiatory affords LDS women the unusual experience of exercising what appears to be a kind of priesthood authority. However, while most LDS rituals require the officiator to state that he is acting by the authority of the Melchizedek priesthood, in the case of the initiatory the officiator does not specify by what authority he or she acts, merely that he or she has authority.''
					}
			I wash you preparatory to your receiving your anointings [for and in behalf of \rule{1cm}{0.15mm}, who is dead], that you may become clean from the blood and sins of this generation.
			
			\textit{[While pronouncing the blessings which follow, the officiator touches each part of the body as it is named.]}
			
			\noindent I wash your head, that your brain and your intellect may be clear and active;\\
			your ears, that you may hear the word of the Lord;\\
			your eyes, that you may see clearly and discern between truth and error;\\
			your nose, that you may smell;\\
			your lips, that you may never speak guile;\\
			your neck, that it may bear up your head properly;\\
			your shoulders, that they may bear the burdens that shall be placed thereon;\\
			your back, that there may be marrow in the bones and in the spine;\\
			your breast, that it may be the receptacle of pure and virtuous principles;\\
			your vitals and bowels, that they may be healthy and perform their proper functions;\\
			your arms and hands, that they may be strong and wield the sword of justice in defense of truth and virtue;\\
			your loins, that you may be fruitful and multiply and replenish the earth, that you might have joy in your posterity;\\
			your legs and feet, that you might run and not be weary, and walk and not faint.
			
			\textit{[A second officiator enters. Both officiators place their hands on the initiate's head, and the second officiator seals the washing as follows.]}
			
			Brother \rule{1cm}{0.15mm}, having authority, we lay our hands upon your head [for and in behalf of \rule{1cm}{0.15mm}, who is dead] and seal upon you this washing, that you may become clean from the blood and sins of this generation through your faithfulness;\\ in the name of Jesus Christ. Amen.
			
		% -----
		\subsubsection{Anointing}
		% -----
			\label{TEMPLE-ENDOWMENT-INIT-PRE-2005-ANOINT}
		
			\textit{[An officiator places oil on the initiate's head while pronouncing the following words.]}
			
			Brother \rule{1cm}{0.15mm}, having authority, I pour this holy anointing oil upon your head%
				\footnote{%
					% \label{}%
					See the scriptures above, where you get an initiatory-like ordinance in the Bible---in particular I find the \hyperref[MATTHEW-26-6]{Matthew 26:6--13} ones interesting, because Christ says, ``For in that she hath poured this ointment on my body, she did it afor my burial.'' The word for ``ointment'' there is \gr{μύρον}, which is not oil but ``ointment, perfume'' (BDAG 586a), used to embalm bodies (\hyperref[LUKE-23-56]{Luke 23:56}). See also \hyperref[MARK-14-3]{Mark 14:3--9}.
					}
			[for and in behalf of \rule{1cm}{0.15mm}, who is dead] and anoint you preparatory to your becoming a king and a priest unto the most high God, hereafter to rule and reign in the house of Israel forever.
			
			\textit{[While pronouncing the blessings which follow, the officiator touches each part of the body as it is named.]}
			
			\noindent I anoint your head, that your brain and your intellect may be clear and active;\\
			your ears, that you may hear the word of the Lord;\\
			your eyes, that you may see clearly and discern between truth and error;\\
			your nose, that you may smell;\\
			your lips, that you may never speak guile;\\
			your neck, that it may bear up your head properly;\\
			your shoulders, that they may bear the burdens that shall be placed thereon;\\
			your back, that there may be marrow in the bones and in the spine;\\
			your breast, that it may be the receptacle of pure and virtuous principles;\\
			your vitals and bowels, that they may be healthy and perform their proper functions;\\
			your arms and hands, that they may be strong and wield the sword of justice
			in defense of truth and virtue;\\
			your loins, that you may be fruitful and multiply and replenish the earth,
			that you might have joy in your posterity;\\
			your legs and feet, that you might run and not be weary, and walk and not faint.
			
			\textit{[A second officiator enters. Both officiators place their hands on the initiate's head, and the second officiator confirms the anointing as follows.]}
			
			Brother \rule{1cm}{0.15mm}, having authority, we lay our hands upon your head [for and in behalf of \rule{1cm}{0.15mm}, who is dead] and confirm upon you this anointing, wherewith you have been anointed in the temple of our God preparatory to becoming a king and a priest unto the most high God, hereafter to rule and reign in the house of Israel forever, and seal upon you all the blessings hereunto appertaining, through your faithfulness; in the name of Jesus Christ. Amen.
			
		% -----
		\subsubsection{The Garment}
		% -----
			\label{TEMPLE-ENDOWMENT-INIT-PRE-2005-GARMENT}
		
			\textit{[An officiator clothes the initiate in the garment. The officiator then pronounces the following words.]}
			
			Brother \rule{1cm}{0.15mm}, having authority, I place this garment upon you [for and in behalf of \rule{1cm}{0.15mm}, who is dead], which you must wear throughout your life. It represents the garment given to Adam when he was found naked in the garden of Eden and is called the garment of the holy priesthood.
			
			Inasmuch as you do not defile it, but are true and faithful to your covenants, it will be a shield and a protection to you against the power of the destroyer until you have finished your work on the earth.
			
		% -----
		\subsubsection{The New Name}
		% -----
			\label{TEMPLE-ENDOWMENT-INIT-PRE-2005-NAME}
		
			\textit{[In the case of a living endowment, the officiator who clothed the initiate in the garment continues as follows.]}
			
			With this garment, I give you a new name, which you should always remember and which you must keep sacred and never reveal, except at a certain place that will be shown you hereafter.
			
			The name is \rule{1cm}{0.15mm}.
			
			\textit{[In the case of a person receiving the endowment on behalf of the dead, the person is presented to an officiator, who pronounces the following words.]}
			
			Brother \rule{1cm}{0.15mm}, having authority, I give you a new name for and in behalf of \rule{1cm}{0.15mm}, who is dead, which you should always remember and which you must keep sacred and never reveal, except at a certain place that will be shown you hereafter.

			The name is \rule{1cm}{0.15mm}.